
Quantifying and tackling uncertainty in deep learning is one of the most challenging problems, mainly due to the inherent randomness of the real world and the presence of noisy data. Bayesian inference provides a robust framework for understanding complex data, allowing for probabilistic interpretation of deep learning models and reasoning under uncertainty. This approach not only facilitates predictions but also enables the quantification of uncertainty. A primary challenge in this domain is the computation and sampling from intricate distributions, mainly when dealing with deep learning models. One effective strategy to tackle this issue is variational inference, which seeks to approximate the true posterior distribution with simpler forms, known as approximate posteriors, while optimizing a variational lower bound. Several techniques have been developed in this area, including those by \citet{kingma2013auto, kingma2015variational}, and \citet{blundell2015weight}, who extended the Gaussian variational posterior approximation for neural networks, as well as \citet{gupta2018matrix}, who enhanced the flexibility of posterior approximations. In addition to variational methods, various particle sampling techniques have been proposed for Bayesian inference, especially in scenarios requiring multiple models. Notable particle sampling methods include Hamiltonian Monte Carlo (HMC) \citep{Neal_Bayesian}, Stochastic Gradient Langevin Dynamics (SGLD) \citep{sgld}, Stochastic Gradient HMC (SGHMC) \citep{chen2014stochastic}, and Stein Variational Gradient Descent (SVGD) \citep{svgd}. Each method contributes to a deeper understanding and more practical application of Bayesian inference in deep learning.

Besides quantifying uncertainty, tackling overfitting is a significant challenge in machine learning. Overfitting often occurs when the training process gets stuck in local minima, leading to a model that fails to generalize well to unseen data. This problem is mainly due to loss functions' high-dimensional and non-convex nature, which often exhibit multiple local minima in the loss landscape. In standard deep network training, flat minimizers improve model generalization \citep{keskar2017large, kaddour, li2022}. Among the flat minimizers, Sharpness-Aware Minimization (SAM) \citep{sam} has emerged as a practical approach by concurrently minimizing the empirical loss and reducing the sharpness of the loss function. Recently, SAM has demonstrated its versatility and effectiveness across a wide range of tasks, including meta-learning \citep{abbas2022sharp}, vision models \citep{chen2021vision}, and language models \citep{bahri-etal-2022-sharpness}.
\begin{figure}[t]
    \centering
    \includegraphics[width=0.45\textwidth]{figs/schematic.pdf}
    \caption{Schematic of SAM w. independent particles {\color{mygreen}(green)}, SVGD {\color{myorange}(orange)}, and our FHBI (black) updates. SAM's particles are not aware of others' trajectories. SVGD only seeks the modes and promotes \textit{spatial} diversity. FHBI seeks the modes, minimizes sharpness, and promotes \textit{spatial} and \textit{angular} diversity.}
    \label{fig:schematic}
\end{figure}

\paragraph{Contribution.} We combine flat minimizers and particle-based Bayesian methods to introduce a novel Bayesian inference framework with improved generalization ability. To accomplish this, we first present Theorem \ref{theorem: flow sharpness}, which strengthens prior generalization bounds from \textit{finite-dimensional} Euclidean spaces to the reproducing kernel Hilbert spaces (RKHS), which are broader functional spaces and are typically \textit{infinite-dimensional}. Notably, this theorem introduces the notion of \textit{functional sharpness} that offers an insight to improve the generalization ability of current particle-sampling methods. Subsequently, Theorem \ref{theorem: pac-bayes} translates these notions of functional sharpness and generalization from the context of RKHS to the context of Bayesian inference. This analysis establishes a connection between the population and empirical KL loss, providing a strategy to enhance generalization by minimizing the population KL loss. Motivated by these two theorems, we present Flat Hilbert Bayesian Inference (FBVI), a practical algorithm that employs a dual-step \textit{functional sharpness-aware} update procedure in RKHS. This approach improves the generalization of sampled particles, thereby enhancing the quality of the ensemble. Overall, our contributions are as follows:
\begin{enumerate}
\item We present a theoretical analysis that characterizes generalization ability over the functional space. This analysis generalizes prior works from the Euclidean space to infinite-dimensional functional space, thereby introducing the notion of \textit{functional sharpness} i.e., the sharpness of the functional spaces.
\item Building on this theoretical foundation, we propose a practical particle-sampling algorithm that enhances the generalization ability over existing methods. We conducted extensive experiments comparing our Flat Hilbert Bayesian Inference (FHBI) algorithm with nine baselines on the \texttt{VTAB-1K} benchmark, which includes 19 datasets across various domains and semantics. Experimental results demonstrated that our algorithm outperforms these baselines by notable margins.
\end{enumerate}
%\textbf{Organization.} {\color{blue} State the organization of the paper.}

\paragraph{Organization.} The paper is structured as follows: Section \ref{section: related works} reviews the related works on Bayesian inference and the development of flat minimizers. Section \ref{section: preliminaries} provides the necessary background and notations. Section \ref{section: methodology} discusses the motivation and theoretical development behind our sharpness-aware particle-sampling approach. Section \ref{section: experiments} presents experimental results, comparing our algorithm against various Bayesian inference baselines across diverse settings. Section \ref{section: ablation} offers a deeper analysis of FHBI's behavior to gain further insight into its effectiveness over the baselines.


% \begin{wrapfigure}{r}{0.5\textwidth}
%     \vspace*{-\baselineskip}
%   \begin{center}
%     \includegraphics[width=0.5\textwidth]{figs/description.pdf}
%   \end{center}
%   \caption{Schematic of FHBI update. Besides driving the particles to the modes, FHBI also encourages diverse travelling directions and pushes the particles away from each other.}
%   \label{fig: architecture}
%   \vspace*{-\baselineskip}
% \end{wrapfigure}

